\documentclass[11pt,letterpaper,oneside]{report}
\usepackage[T1]{fontenc}
\usepackage[utf8]{inputenc}
\usepackage{lmodern}
\usepackage[letterpaper,margin=0.92in,headheight=15pt]{geometry}
\usepackage{amsmath,amssymb,amsthm,mathtools}
\usepackage{microtype}
\usepackage[dvipsnames]{xcolor}
\usepackage{graphicx}
\usepackage{booktabs,longtable,array,tabularx}
\usepackage{enumitem}
\usepackage{fancyhdr}
\usepackage{titlesec}
\usepackage[most]{tcolorbox}
\usepackage{listings}
\usepackage{tikz}
\usetikzlibrary{positioning}
\usepackage{hyperref}
\hypersetup{colorlinks=true,linkcolor=MidnightBlue,urlcolor=RoyalBlue,pdfauthor={Independent Mathematics and Engineering Study Suite},pdftitle={Mathematical Modeling and Engineering Applications Handbook}}
\definecolor{ink}{HTML}{162639}
\definecolor{navy}{HTML}{173B61}
\definecolor{blue}{HTML}{2874A6}
\definecolor{teal}{HTML}{007C83}
\definecolor{amber}{HTML}{8B5E16}
\definecolor{pale}{HTML}{EDF4F8}
\definecolor{paleteal}{HTML}{ECF8F5}
\definecolor{paleamber}{HTML}{FFF6E8}
\titleformat{\chapter}[display]{\normalfont\huge\bfseries\color{navy}}{\filleft\fontsize{43}{46}\selectfont\thechapter}{0.4ex}{\titlerule\vspace{1ex}}[\vspace{1ex}\titlerule]
\titleformat{\section}{\normalfont\Large\bfseries\color{navy}}{\thesection}{0.8em}{}
\titleformat{\subsection}{\normalfont\large\bfseries\color{teal}}{\thesubsection}{0.75em}{}
\pagestyle{fancy}\fancyhf{}\fancyhead[L]{\small\color{navy}Differential Equations Engineering Suite}\fancyhead[R]{\small\color{navy}Volume 7 of 8}\fancyfoot[C]{\thepage}\renewcommand{\headrulewidth}{0.3pt}
\setlist[itemize]{leftmargin=*,itemsep=3pt,topsep=5pt}\setlist[enumerate]{leftmargin=*,itemsep=4pt,topsep=5pt}
\setlength{\parskip}{5pt}\setlength{\parindent}{0pt}\setlength{\emergencystretch}{2em}
\newtheorem{theorem}{Theorem}[chapter]
\newtheorem{proposition}[theorem]{Proposition}
\newtheorem{lemma}[theorem]{Lemma}
\theoremstyle{definition}\newtheorem{definition}[theorem]{Definition}
\newtheorem{example}[theorem]{Worked example}
\theoremstyle{remark}\newtheorem{remark}[theorem]{Remark}
\newtcolorbox{bookground}[1][]{enhanced,breakable,colback=pale,colframe=blue,boxrule=0.6pt,title={TEXTBOOK ALIGNMENT},fonttitle=\sffamily\bfseries\small,#1}
\newtcolorbox{modern}[1][]{enhanced,breakable,colback=paleteal,colframe=teal,boxrule=0.6pt,title={MODERN COMPUTATIONAL EXTENSION},fonttitle=\sffamily\bfseries\small,#1}
\newtcolorbox{caution}[1][]{enhanced,breakable,colback=paleamber,colframe=amber,boxrule=0.6pt,title={ASSUMPTION / VALIDATION},fonttitle=\sffamily\bfseries\small,#1}
\newcommand{\dd}{\,\mathrm{d}}
\newcommand{\R}{\mathbb{R}}
\newcommand{\lap}{\mathcal{L}}
\newcommand{\norm}[1]{\left\lVert#1\right\rVert}
\newcommand{\abs}[1]{\left\lvert#1\right\rvert}
\lstdefinestyle{engpy}{language=Python,basicstyle=\ttfamily\footnotesize,keywordstyle=\color{RoyalBlue}\bfseries,commentstyle=\color{gray},stringstyle=\color{teal},backgroundcolor=\color{pale},frame=single,rulecolor=\color{blue!30},breaklines=true,columns=fullflexible,keepspaces=true,showstringspaces=false,tabsize=4,upquote=true}
\lstset{style=engpy}
\begin{document}
\begin{titlepage}
\pagecolor{navy}\color{white}
\vspace*{0.75in}
{\sffamily\bfseries\large MATHEMATICS \& ENGINEERING DOCUMENTATION SUITE\\[12pt]}
{\color{cyan!40}\rule{\textwidth}{2pt}}\\[1.0in]
{\sffamily\bfseries\fontsize{29}{36}\selectfont Mathematical Modeling and Engineering Applications Handbook\par}
\vspace{0.45in}
{\Large Volume 7 of 8}\par\vspace{0.55in}
{\color{cyan!30}\rule{0.50\textwidth}{1pt}}\par\vspace{0.3in}
{\large Independent technical companion aligned with the\par}
{\large 13-chapter organization of Nagle, Saff \& Snider (7th ed.)\par}
\vfill
{\large Original derivations, worked examples, verification and engineering practice\par}\vspace{0.3in}
{\small Prepared October 2026 \quad\textbullet\quad pdf\LaTeX{} compatible\par}
\end{titlepage}
\nopagecolor\color{ink}\pagenumbering{roman}
\chapter*{Reader's guide and source boundaries}\addcontentsline{toc}{chapter}{Reader's guide and source boundaries}
This independent study resource is organized around \emph{Fundamentals of Differential Equations and Boundary Value Problems}, Seventh Edition, by R.~Kent Nagle, Edward~B. Saff, and Arthur~David Snider (Pearson, 2018; ISBN 978-0-321-97710-6). The book's chapter and section numbers are given as \emph{navigation references}, not as claims of reproduced text or proofs.

\begin{bookground}
\textbf{Textbook-grounded scope:} topics, terminology, and chapter/section organization corresponding to the cited textbook. All explanations, derivations, exercise prompts, and diagrams in this volume are newly authored and may emphasize different examples. Check the textbook itself for its exact hypotheses, proofs, exercises, or nomenclature in a particular section.
\end{bookground}
\begin{modern}
\textbf{Distinct extension layer:} Python/NumPy/SciPy workflows, engineering QA and reproducibility practices, automated convergence studies, diagnostic plots, and implementation guidance. These are \emph{not} represented as content supplied by the textbook. API calls are examples and should be verified against the installed package version.
\end{modern}
\begin{caution}
\textbf{Mathematical qualification:} formulas require their stated regularity, domain and boundary/initial conditions. Numerical agreement does not replace a proof; computational outputs are not certified error bounds unless explicitly established.
\end{caution}
\tableofcontents\clearpage\pagenumbering{arabic}
\chapter{Engineering modeling from first principles}
\begin{bookground}Primary textbook anchors: Chapter 3 (Mathematical Models and Numerical Methods Involving First-Order Equations), Chapter 4 vibrations, Chapter 5 systems, Chapter 10 PDEs, and Chapter 12 autonomous stability. The engineering lifecycle and modern data-science practices are newly authored additions.\end{bookground}
\section{Model development workflow}
\begin{enumerate}
\item Define the prediction question, observable outputs and operating envelope.
\item Draw a control volume or free-body diagram and identify state variables.
\item Apply conservation, kinematics or circuit laws, with documented constitutive relations.
\item Separate assumptions from laws (well-mixed, linear, incompressible, constant coefficients, etc.).
\item Specify initial/boundary conditions, disturbances and parameter units.
\item Establish analytic limiting cases and calculate dimensionless groups.
\item Select an analytical or numerical solution, verify it, and validate with data.
\end{enumerate}
A mathematically solvable differential equation is not necessarily a validated physical model. Conversely, a physically accurate model may require numerical methods rather than a closed form.
\section{Dimensional consistency and scaling}
For any valid sum, terms must share dimensions. If $C\dot T=-H(T-T_a)+Q$, then $C$ has dimensions energy per temperature, $H$ power per temperature and $Q$ power. The time constant $\tau=C/H$ has units of time. Use such consistency checks before parameter estimation or simulation.
\chapter{Compartmental analysis and mixing}
\begin{bookground}Correspondence: \S3.2 Compartmental Analysis; \S5.1 Interconnected Fluid Tanks.\end{bookground}
\section{One perfectly mixed vessel}
Let $V$ be constant liquid volume, $q$ inflow/outflow volume per unit time, $c_{\rm in}(t)$ inlet concentration and $c(t)$ internal concentration. Conservation gives
\[\frac{d}{dt}(Vc)=q c_{\rm in}-qc,\qquad c'+\frac qV c=\frac qV c_{\rm in}.\]
For constant inlet concentration $c_\infty$ and initial $c_0$,
\[c(t)=c_\infty+(c_0-c_\infty)e^{-qt/V}.\]
The time constant is $V/q$. A well-mixed assumption predicts uniform concentration; stratification, short-circuiting and reactive chemistry require extensions.
\begin{example}[Residence time]
For $V=1000\,\mathrm{L}$ and $q=100\,\mathrm{L/min}$, the turnover time is $10$ minutes. A concentration difference from its steady state decays by a factor $e^{-1}$ after 10 minutes in the idealized model.
\end{example}
\section{Interconnected tanks}
For two equal-volume tanks with a bidirectional exchange rate $q$, concentrations satisfy $c_1'=\alpha(c_2-c_1)$ and $c_2'=\alpha(c_1-c_2)$ where $\alpha=q/V$. Sum is constant while the difference $d=c_1-c_2$ satisfies $d'=-2\alpha d$. Thus differences decay exponentially without changing total dissolved mass. This offers a direct analytic test for a system integrator.
\chapter{Thermal dynamics of buildings}
\begin{bookground}Correspondence: \S3.3 Heating and Cooling of Buildings; related heat transfer in Chapter 10.\end{bookground}
\section{Lumped RC thermal model}
Let indoor temperature $T_i$, outdoor temperature $T_o(t)$, thermal capacitance $C$, envelope conductance $H$, and controlled heat input $Q(t)$. Energy balance gives
\[C\frac{dT_i}{dt}=H[T_o(t)-T_i(t)]+Q(t).\]
For constant $T_o,Q$ define $T_*=T_o+Q/H$ and $\tau=C/H$. Then $T_i(t)=T_*+[T_i(0)-T_*]e^{-t/\tau}$. This assumes a single well-mixed thermal zone and constant parameters.
\section{Engineering limitations}
Real buildings contain multiple thermal masses, solar gains, infiltration, humidity coupling, occupant disturbances and actuator constraints. Upgrade to a multi-zone state-space model where warranted, rather than interpreting an inaccurate single-time-constant fit as a numerical error. A temperature sensor measures a local environment, which may not equal the state assumed by a one-node model.
\section{Thermal model QA}
Check bounds under bounded heating and ambient conditions, correct response direction to turning a heater on/off, time-constant fit to a step change, and physical units. If a thermostat causes on/off switching, model the control as a hybrid forcing and resolve switching times accurately.
\chapter{Newtonian mechanics and vibration}
\begin{bookground}Correspondence: \S3.4 Newtonian Mechanics and \S4.1, \S4.9--4.10 mass-spring oscillators and mechanical vibrations; coupled oscillators in \S5.6.\end{bookground}
\section{Force balance}
Newton's second law in one dimension is $m x''=\sum F$. For linear elasticity, viscous damping and forcing,
\[m x''+c x'+k x=F(t),\qquad m>0,\ k>0,\ c\ge0.\]
The undamped natural angular frequency is $\omega_n=\sqrt{k/m}$ and damping ratio $\zeta=c/(2\sqrt{mk})$. Cases $\zeta<1$, $\zeta=1$, and $\zeta>1$ correspond to underdamped, critically damped and overdamped homogeneous responses.
\section{Step response and energy}
For constant force $F_0$, the new static equilibrium is $x_*=F_0/k$. Let $z=x-x_*$; then $m z''+cz'+kz=0$. For $F(t)=0$ mechanical energy is $E=\frac12m(x')^2+\frac12kx^2$ and $\dot E=-c(x')^2$. Monitor the latter invariant to check signs in a numerical implementation.
\begin{example}[Natural frequency]
For $m=2\,\mathrm{kg}$ and $k=50\,\mathrm{N/m}$, $\omega_n=5\,\mathrm{rad/s}$ and $f_n=5/(2\pi)\,\mathrm{Hz}\approx0.796\,\mathrm{Hz}$. If $c=4\,\mathrm{N\,s/m}$, $\zeta=4/(2\sqrt{100})=0.2$.
\end{example}
\section{Harmonic forced response}
Under $F(t)=F_0\cos\omega t$ and $c>0$, a steady-state complex amplitude is
\[\hat x=\frac{F_0}{k-m\omega^2+i c\omega},\quad |\hat x|=\frac{F_0}{\sqrt{(k-m\omega^2)^2+(c\omega)^2}}.\]
The observed displacement is the real part of $\hat x e^{i\omega t}$ plus a decaying transient. For nonzero damping the maximum response frequency need not be exactly $\omega_n$.
\chapter{Electrical circuits and control}
\begin{bookground}Correspondence: \S3.5 Electrical Circuits and \S5.7 Electrical Systems; frequency-response modeling project in Chapter 7.\end{bookground}
\section{Series RLC circuit}
For a series RLC branch, Kirchhoff's voltage law is $L\ddot q+R\dot q+q/C=v_s(t)$ with $i=\dot q$. Comparing with mechanical motion reveals $m\leftrightarrow L$, $c\leftrightarrow R$, and $k\leftrightarrow1/C$. The circuit is stable for passive $R>0$, $L,C>0$ under the ideal linear model.
\section{Transfer-function viewpoint}
For zero initial conditions, $I(s)/V_s(s)=s/[Ls^2+Rs+1/C]$. Initial energy requires extra transform terms when present. Pole locations govern homogeneous modes; zeros affect input-output response. Transfer functions characterize linear time-invariant models, not arbitrary nonlinear switched systems.
\section{Component tolerances}
Capacitors and inductors have uncertainty, loss and saturation effects. Evaluate sensitivity to $R,L,C$ rather than quoting response parameters to unjustified precision. Numerical simulation should preserve circuit sign conventions and check voltage/current units throughout.
\chapter{Biomathematics and epidemiological dynamics}
\begin{bookground}Correspondence: \S5.5 Applications to Biomathematics (epidemic and tumor-growth models) and Chapter 12 population/neuronal examples. These illustrations are mathematical exercises, not clinical or public-health forecasts.\end{bookground}
\section{SIR compartment equations}
For a closed population $N=S+I+R$ with infection coefficient $\beta$ and removal coefficient $\gamma$, a standard idealized model is
\[S'=-\beta SI/N,\quad I'=\beta SI/N-\gamma I,\quad R'=\gamma I.\]
Total population remains constant: $S'+I'+R'=0$. Initial infection growth occurs if $\beta S(0)/(\gamma N)>1$ when $I(0)>0$. The simple reproduction ratio $R_0=\beta/\gamma$ assumes a wholly susceptible population and homogeneous mixing.
\section{Boundaries and positivity}
For nonnegative initial conditions and positive parameters, the SIR vector field preserves nonnegative states in the idealized model. Negative numerical compartments indicate a numerical or implementation issue. Contact heterogeneity, stochasticity, demography, spatial effects and parameter changes invalidate naive interpretation of a constant-parameter SIR curve.
\section{Population competition}
The Lotka--Volterra competition family $x'=r_1x(1-(x+\alpha y)/K_1)$ and $y'=r_2y(1-(y+\beta x)/K_2)$ illustrates equilibrium, nullcline and Jacobian analysis. Coexistence depends on interspecific coefficients and capacity parameters. A stability finding is model-relative, not a universally valid ecological prediction.
\chapter{Heat transport and spatial modeling}
\begin{bookground}Correspondence: Chapter 10 heat, wave and Laplace equations, and Chapter 11 eigenvalue BVPs.\end{bookground}
\section{Heat conduction PDE}
A constant-property one-dimensional solid with heat source $Q(x,t)$ satisfies $\rho c_p T_t=k T_{xx}+Q$. Under fixed zero temperature difference at both endpoints and no source, a Fourier sine series describes decay. Thermal diffusivity $\kappa=k/(\rho c_p)$ has units $\mathrm{m^2/s}$.
\section{Characteristic time scale}
The diffusion scale for a length $L$ is $t_d=L^2/\kappa$. Doubling spatial length quadruples the diffusive time scale, holding the material properties fixed. This dimensionless insight helps decide observation horizons and time-step constraints in numerical studies.
\section{When a PDE is necessary}
A lumped temperature ODE neglects internal gradients. When spatial gradient effects drive performance or safety, use an appropriate spatial model and boundary heat-transfer condition, e.g. Robin form $-k\partial_n T=h(T-T_\infty)$ under a consistent outward-normal sign convention.
\chapter{Computational extension: comparative engineering simulations}
\begin{modern}The following implementation uses SciPy \texttt{solve\_ivp} for educational reproducibility. It is not claimed to be textbook code.\end{modern}
\begin{lstlisting}
import numpy as np
from scipy.integrate import solve_ivp

# Damped oscillator: state [position (m), velocity (m/s)]
m, c, k = 2.0, 4.0, 50.0

def oscillator(t, s):
    x, v = s
    return [v, -(c/m)*v - (k/m)*x]

sol = solve_ivp(oscillator, (0, 10), [0.1, 0.0],
                rtol=1e-9, atol=[1e-11, 1e-11],
                dense_output=True)
assert sol.success
xx, vv = sol.y
energy = .5*m*vv**2 + .5*k*xx**2
assert np.max(np.diff(energy)) < 1e-7  # sampled energy

# Closed SIR example: counts or population fractions
beta, gamma = 0.3, 0.1

def sir(t, x):
    S, I, R = x
    N = S + I + R
    return [-beta*S*I/N, beta*S*I/N-gamma*I, gamma*I]

sir_sol = solve_ivp(sir, (0, 80), [990.0, 10.0, 0.0],
                    rtol=1e-9, atol=1e-8)
assert sir_sol.success
assert np.max(np.abs(sir_sol.y.sum(axis=0)-1000)) < 1e-4
\end{lstlisting}
\section{Calibration and uncertainty}
Estimate parameters by minimizing differences between predicted observations and independent data, but guard against \emph{structural nonidentifiability} (different parameters create the same outputs) and \emph{practical nonidentifiability} (data are too weak/noisy to distinguish estimates). Report sensor calibration, observation timing and fitted noise assumptions.
\section{Scenario and sensitivity analysis}
Evaluate multiple plausible parameter configurations and forcing scenarios. A local normalized sensitivity can be estimated as $S_p\approx(p/y)[y(p+\Delta p)-y(p-\Delta p)]/(2\Delta p)$ for a nonzero output $y$ and parameter $p$; disclose the perturbation size and output metric. Near discontinuities or bifurcations, a local derivative may not capture the full risk.
\chapter{Engineering capstones and acceptance criteria}
\section{Capstone A: building heating controller}
Develop a one-zone RC thermal model; calculate a time constant; incorporate a heater on/off input; estimate $C$ and $H$ from synthetic or measured temperature data; quantify parameter uncertainty. \emph{Acceptance:} unit consistency, verified steady state, solver status, and error plots on withheld observations.
\section{Capstone B: coupled oscillator design}
Construct a two-degree-of-freedom mass-spring system. Verify positive definiteness of mass and stiffness matrices, compute natural frequencies, simulate free response, and compare modal reconstruction with direct integration. \emph{Acceptance:} mode orthogonality in the mass inner product and energy conservation in the undamped limit.
\section{Capstone C: mixed-tank incident response}
Extend the two-tank model with changing inflow concentrations and independent inflow/outflow rates. Check positivity and mass balance. \emph{Acceptance:} well-specified mixing assumptions, boundary of validity and separate numerical/model error estimates.
\section{Capstone D: thermal diffusion}
Solve a one-dimensional rod heating/cooling problem using sine modes and method of lines. Refine spatial and temporal grids and compare thermal energy dissipation. \emph{Acceptance:} correct boundary data, documented convergence and reproducible plots.
\section{Capstone E: infectious disease idealization}
Analyze the qualitative thresholds of an SIR system under several initial susceptible fractions. Do not claim real-world forecasting accuracy without empirical calibration. \emph{Acceptance:} nonnegative states, conserved total population, explicit parameter interpretation and limitations.

\appendix
\chapter{References and reproducibility}
\begin{enumerate}
\item R.~K. Nagle, E.~B. Saff, and A.~D. Snider, \emph{Fundamentals of Differential Equations and Boundary Value Problems}, 7th ed., Pearson, 2018, ISBN 978-0-321-97710-6. Primary source for the \emph{topic map and terminology}; cited locally by chapter and section.
\item SciPy project, \emph{solve\_ivp} reference, \url{https://docs.scipy.org/doc/scipy/reference/generated/scipy.integrate.solve_ivp.html} (modern extension).
\item SciPy project, \emph{solve\_bvp} reference, \url{https://docs.scipy.org/doc/scipy/reference/generated/scipy.integrate.solve_bvp.html} (modern extension).
\item SciPy project, \emph{eigsh} reference, \url{https://docs.scipy.org/doc/scipy/reference/generated/scipy.sparse.linalg.eigsh.html} (modern extension).
\item Python, NumPy, SciPy, and SymPy project documentation; computational checks require recording package versions and tolerances.
\end{enumerate}
\section*{Verification record template}
\begin{tabularx}{\textwidth}{>{\bfseries}p{0.30\textwidth}X}\toprule
Problem identifier & \rule{0pt}{15pt}\ \\ \midrule
ODE/PDE and domain & \\[8pt]
Initial/boundary conditions & \\[8pt]
Units and parameter values & \\[8pt]
Analytic residual & \\ [8pt]
Grid/step refinement & \\ [8pt]
Independent solver comparison & \\ [8pt]
Version and reproducibility notes & \\ [8pt]\bottomrule
\end{tabularx}
\end{document}
